\chapter{Alternativas de pesquisa}\label{comparauxe7uxe3o-entre-possuxedveis-projetos}

{
\begingroup
\small
\singlespacing
\setlength{\tabcolsep}{4pt}
\begin{longtable}[]{@{}
  >{\raggedright\arraybackslash}p{(\linewidth - 4\tabcolsep) * \real{0.3333}}
  >{\raggedright\arraybackslash}p{(\linewidth - 4\tabcolsep) * \real{0.3333}}
  >{\raggedright\arraybackslash}p{(\linewidth - 4\tabcolsep) * \real{0.3333}}@{}}
\toprule\noalign{}
\begin{minipage}[b]{\linewidth}\raggedright
Linha
\end{minipage} & \begin{minipage}[b]{\linewidth}\raggedright
Implementação e validação
\end{minipage} & \begin{minipage}[b]{\linewidth}\raggedright
Avaliação para 24 meses
\end{minipage} \\ \addlinespace[4pt]
\midrule\noalign{}
\endhead
\bottomrule\noalign{}
\endlastfoot
\textbf{Inferência de massa com PTAs: compressão em frequência e vieses}
& Simulação, resposta dispersiva, comparação entre métodos e calibração
por injeções. & \textbf{Recomendada.} Pergunta estreita, referências de
2026 e resultado útil mesmo sem nova detecção. \\ \addlinespace[4pt]
Polarizações exatas e consistência de teorias massivas & Álgebra
simbólica, vínculos e comparação com Fierz--Pauli. & Excelente primeiro
capítulo. Para artigo próprio precisa superar resultados conhecidos,
especialmente os de 2019. \\ \addlinespace[4pt]
Polarizações em LISA ou LISA--Taiji & Resposta interferométrica, formas
de onda e inferência. & Alternativa se houver orientação especializada.
Há trabalhos extensos de 2026; o recorte deve ser muito específico. \\ \addlinespace[4pt]
Correlações de três pontos para polarizações escalares & Derivação de
resposta e simulação de fundos não gaussianos. & Mais exploratória e
teórica; maior risco para mestrado. Deve ficar como alternativa, não ser
somada ao projeto principal. \\ \addlinespace[4pt]
\end{longtable}
\endgroup
}

Para LISA, os pontos de partida atuais incluem \textcite{akama2026} e
\textcite{mu2026}. Para correlações de três pontos,
\textcite{jimenez2026} propõem um diagnóstico de polarizações escalares
sob hipóteses específicas de média sobre polarizações. Uma extensão para
ondas massivas precisaria verificar dispersão, homogeneidade,
estacionariedade e condições de ressonância; um fundo gaussiano tem
bispectro nulo e não pode ser usado para prometer esse sinal.

A extensão cosmológica de Visser sugerida no TG também tem antecedentes
diretos, inclusive de seus colaboradores: conservação de
energia-momento\footnote{Disponível em:
  \url{https://arxiv.org/abs/0710.1077}. Consulta em 10 set. 2026.} e
aplicação cosmológica\footnote{Disponível em:
  \url{https://arxiv.org/abs/0907.5190}. Consulta em 10 set. 2026.}. Ela
demandaria um problema novo e uma análise de estabilidade mais ampla,
sendo menos indicada como primeira escolha neste caso.
